\section{Derivatives, recurrences, and scaling series}\label{sec:derivatives}

\subsection{Four first-order derivative formulas}

The identity $\frac{\dd}{\dd z}z^s=sz^{s-1}$ combines with a single Gamma
recurrence.  For example,
\[
 (s-a_1+1)\Gamma(1-a_1+s)=\Gamma(2-a_1+s),
\]
which lowers $a_1$ by one.  Depending on the block, the sign is reversed.

\begin{theorem}[First-order formulas]\label{thm:first-derivatives}
On every certified differentiability domain,
\begin{align}
\frac{\dd}{\dd z}\!
 \left[z^{1-a_1}G(\mathbf a;\mathbf b\mid z)\right]
&=z^{-a_1}G(a_1-1,a_2,\ldots;\mathbf b\mid z),&&n\ge1,
\label{eq:d-a1}\\
\frac{\dd}{\dd z}\!
 \left[z^{1-a_p}G(\mathbf a;\mathbf b\mid z)\right]
&=-z^{-a_p}G(a_1,\ldots,a_p-1;\mathbf b\mid z),&&n<p,
\label{eq:d-ap}\\
\frac{\dd}{\dd z}\!
 \left[z^{-b_1}G(\mathbf a;\mathbf b\mid z)\right]
&=-z^{-1-b_1}G(\mathbf a;b_1+1,b_2,\ldots\mid z),&&m\ge1,
\label{eq:d-b1}\\
\frac{\dd}{\dd z}\!
 \left[z^{-b_q}G(\mathbf a;\mathbf b\mid z)\right]
&=z^{-1-b_q}G(\mathbf a;b_1,\ldots,b_q+1\mid z),&&m<q.
\label{eq:d-bq}
\end{align}
\end{theorem}
\begin{proof}
Differentiate a finite contour integral or finite residue partial sum.  In each
case, collect the factor $s+1-a_i$ or $s-b_j$ with the affected Gamma or
reciprocal-Gamma factor and use $\Gamma(w+1)=w\Gamma(w)$.  Uniform derivative
tails follow from \cref{thm:analyticity}.  The sign in
\eqref{eq:d-ap} and \eqref{eq:d-b1} is the sign obtained from
$a_p-s=-(s-a_p)$ or $b_1-s=-(s-b_1)$.
\end{proof}

Expanding \eqref{eq:d-a1} gives the useful Euler-derivative identity
\begin{equation}\label{eq:theta-shift}
 \thetaop G(\mathbf a;\mathbf b\mid z)
 =G(a_1-1,a_2,\ldots;\mathbf b\mid z)+(a_1-1)G(\mathbf a;\mathbf b\mid z).
\end{equation}
The other three analogous identities are obtained in the same finite way.

\subsection{Arbitrary derivative order}

For $h\in\N$, the falling factorial
$s^{\underline h}=s(s-1)\cdots(s-h+1)$ is represented by adding a top parameter
$0$ and a bottom parameter $h$, because
\begin{equation}\label{eq:falling-gamma}
 \frac{\Gamma(1+s)}{\Gamma(1-h+s)}=s^{\underline h}.
\end{equation}

\begin{theorem}[Higher derivatives]\label{thm:higher-derivatives}
For each fixed $h$,
\begin{align}
z^h\frac{\dd^h}{\dd z^h}G_{p,q}^{m,n}(\mathbf a;\mathbf b\mid z)
&=G_{p+1,q+1}^{m,n+1}(0,\mathbf a;\mathbf b,h\mid z)\label{eq:high-d-one}\\
&=(-1)^hG_{p+1,q+1}^{m+1,n}(\mathbf a,0;h,\mathbf b\mid z).
\label{eq:high-d-two}
\end{align}
For the inverse argument,
\begin{align}
z^h\frac{\dd^h}{\dd z^h}G_{p,q}^{m,n}(\mathbf a;\mathbf b\mid z^{-1})
&=G_{p+1,q+1}^{m+1,n}(\mathbf a,1-h;1,\mathbf b\mid z^{-1})\notag\\
&=(-1)^hG_{p+1,q+1}^{m,n+1}(1-h,\mathbf a;\mathbf b,1\mid z^{-1}).
\label{eq:high-d-inverse}
\end{align}
\end{theorem}
\begin{proof}
The $h$th derivative of $z^s$ is $s^{\underline h}z^{s-h}$.
Equation \eqref{eq:falling-gamma} inserts the two new parameters and proves
\eqref{eq:high-d-one}.  The block-transfer identity
\eqref{eq:block-transfer} with difference $h$ proves
\eqref{eq:high-d-two}.  Apply inversion and the finite chain rule to obtain
\eqref{eq:high-d-inverse}.  Since $h$ is externally fixed, all derivative
error bounds remain primitive recursive.
\end{proof}

\subsection{Contiguous recurrences}

Equating two versions of $\thetaop G$ gives the source page's recurrences.
They are therefore not new analytic theorems; they are finite linear
combinations of \cref{thm:first-derivatives}.

\begin{theorem}[Four displayed contiguous relations]\label{thm:contiguous}
Using the shorthand $G=G(\mathbf a;\mathbf b\mid z)$,
\begin{align}
(a_p-a_1)G
&=G(a_1-1,a_2,\ldots;\mathbf b)
 +G(a_1,\ldots,a_p-1;\mathbf b),&&1\le n<p,
\label{eq:contig-aa}\\
(b_1-b_q)G
&=G(\mathbf a;b_1+1,b_2,\ldots)
 +G(\mathbf a;b_1,\ldots,b_q+1),&&1\le m<q,
\label{eq:contig-bb}\\
(b_1-a_1+1)G
&=G(a_1-1,a_2,\ldots;\mathbf b)
 +G(\mathbf a;b_1+1,b_2,\ldots),&&m,n\ge1,
\label{eq:contig-ab}\\
(a_p-b_q-1)G
&=G(a_1,\ldots,a_p-1;\mathbf b)
 +G(\mathbf a;b_1,\ldots,b_q+1),&&n<p,\ m<q.
\label{eq:contig-ba}
\end{align}
\end{theorem}
\begin{proof}
Expand the left derivatives in \eqref{eq:d-a1}--\eqref{eq:d-bq}, solve each for
$\thetaop G$, and subtract the appropriate pair.  This uses only distributivity
and integer/rational coefficient arithmetic, so it compiles to a short object
proof once the four derivative lemmas are available.
\end{proof}

The two omitted diagonal recurrences mentioned by the page follow by adding or
subtracting \eqref{eq:contig-aa}--\eqref{eq:contig-ba}; the checker can generate
them automatically.

\subsection{Scaling/Taylor formulas}

The four source expansions are Taylor series in the scale parameter $w$ around
$w=1$.  We prove one; the others follow from the corresponding first-order
formula.  Put
\[
 H(w)=w^{-b_1}G(\mathbf a;\mathbf b\mid wz).
\]
Using \eqref{eq:d-b1} and the chain rule gives
\[
 H^{(h)}(w)=(-1)^h w^{-b_1-h}
 G(\mathbf a;b_1+h,b_2,\ldots\mid wz),
 \qquad
 \frac{H^{(h)}(1)}{h!}
 =\frac{(-1)^h}{h!}G(\mathbf a;b_1+h,b_2,\ldots\mid z).
\]
Rewriting $(-1)^h(w-1)^h=(1-w)^h$ gives the first displayed
series.  The other three are identical after replacing the first-order shift
identity.

\begin{theorem}[Scaling series]\label{thm:scaling-series}
On every rational disk in the $w$-plane for which the common branch and a
uniform derivative bound are certified,
\begin{align}
G(\mathbf a;\mathbf b\mid wz)
&=w^{b_1}\sum_{h\ge0}\frac{(1-w)^h}{h!}
  G(\mathbf a;b_1+h,b_2,\ldots\mid z),&&m\ge1,
\label{eq:scale-b1}\\
G(\mathbf a;\mathbf b\mid wz)
&=w^{b_q}\sum_{h\ge0}\frac{(w-1)^h}{h!}
  G(\mathbf a;b_1,\ldots,b_q+h\mid z),&&m<q,
\label{eq:scale-bq}\\
G(\mathbf a;\mathbf b\mid z/w)
&=w^{1-a_1}\sum_{h\ge0}\frac{(1-w)^h}{h!}
  G(a_1-h,a_2,\ldots;\mathbf b\mid z),&&n\ge1,
\label{eq:scale-a1}\\
G(\mathbf a;\mathbf b\mid z/w)
&=w^{1-a_p}\sum_{h\ge0}\frac{(w-1)^h}{h!}
  G(a_1,\ldots,a_p-h;\mathbf b\mid z),&&n<p.
\label{eq:scale-ap}
\end{align}
\end{theorem}
\begin{proof}
Taylor's formula with an integral or Cauchy remainder is applied on a certified
rational disk.  The derivatives are replaced by shifted $G$-functions using
\cref{thm:first-derivatives,thm:higher-derivatives}.  A uniform Cauchy bound
$M/R^h$ gives the tail
$M(\abs{w-1}/R)^{N+1}/(1-\abs{w-1}/R)$, with a primitive-recursive cutoff.
This also explains why the page correctly warns that the radius depends on
$z$ and on the function being expanded.
\end{proof}
